\chapter{Introduzione agli assiomi}

La Meccanica Quantistica è basata su tre assimio fondamentali:

\begin{assioma}
    Uno stato fisico $\ket{\psi}$ è un vettore sullo spazio di Hilbert $\mathcal{H}$. Gli stati che consideriamo godono della condizione di normalizzazione 
    \begin{equation}
        \braket{\psi|\psi} = 1
    \end{equation}
    Più propriamente, gli stati sono raggi di modulo $1$ dello spazio $\mathcal{H}$, in quanto sono definiti a meno di una fase globale
    \begin{equation*}
        \ket{\psi} \sim e^{i\alpha} \ket{\psi} \quad , \quad \alpha \in \mathbb{R}
    \end{equation*}
\end{assioma}

\begin{assioma}
    A ciascun osservabile (o variabile dinamica) $O$ è associato un operatore lineare $\hat{O}: \mathcal{H} \to \mathcal{H}$ hermitiano, $\hat{O}^{\dagger} = \hat{O}$. I valori ammessi come esito di una misura sono gli autovalori di $\hat{O}$.
\end{assioma}

\begin{assioma}
    Il valor medio di un osservabile $O$ su un ensamble di sistemi preparati identicamente, cioè il valore atteso dell'osservabile $O$ su uno stato $\ket{\psi}$ è dato dal seguente elemento di matrice 
    \begin{equation}
        \langle O \rangle \equiv \sum_n \lambda_n P_n = \Braket{\psi |\hat{O}|\psi}
    \end{equation}
    dove $\{\lambda_n\}$ sono gli autovalori (quindi assunti discreti) dell'operatori $\hat{O}$, mentre $P_n$ rappresenta la probabilità che l'osservabile $\hat{O}$ assuma il valore $\lambda_n$ quando il sistema è nello stato $\ket{\psi}$. 
\end{assioma}

Tuttavia se non viene data una giustificazione esplicita dell'introduzione di questi assiomi, questi potrebbero sembrare al quanto misteriosi. Di seguito andremo a mostrare come l'introduzione degli assiomi sia giustificata dall'esperienza pratica tramite un esperimento svolto nel $1922$ (quindi prima dell'avvento formale della meccanica quantistica) da Otto Stern e Walther Gerlach.  

\section{Stern-Gerlach}

Otto Stern e Walther Gerlach escogitarono un esperimento per determinare se le particelle avessero un momento angolare intrinseco. 

Supponiamo di considerare particelle neutre e supponiamo che esse abbiamo momento magnetico intriseco $\mathbf{\mu}$. 
Se la particella attraversa un campo magnetico non omogeneo, la forza a un'estremità del dipolo sarà leggermente maggiore di quella all'estremità opposta; questo causa la deflessione della particella. Si ha infatti che 
%
\begin{equation}
    E_m(\mathbf{r}) = - \mathbf{\mu} \cdot \mathbf{B}(\mathbf{r}) \implies \mathbf{F} = - \nabla E_m(\mathbf{r}) = \nabla (\mathbf{\mu} \cdot \mathbf{B}(\mathbf{r}))
\end{equation}
Consideriamo un campo diretto prevalentemente lungo un asse, che scegliamo
essere quello delle $z$, e disomogeneo lungo lo stesso asse. Resta
%
\begin{equation}
    F_z \simeq \mu_z \frac{\partial B_z}{\partial z}
\end{equation}
%
Quindi le particelle sono deflesse verso l'alto o verso il basso a seconda del
segno di $\mu_z$.
Andiamo adesso a costruire uno strumento in grado di generare un campo magnetico fortemente disomogeneo nella regione attraversata dal fascio.\\
L'apparato si compone di quattro elementi in sequenza:
%
\begin{enumerate}
    \item un \textbf{forno}, in cui si vaporizza il materiale (nell'esperimento originale argento) e da cui esce un fascio di atomi neutri;
    \item un sistema di \textbf{fenditure collimatrici}, che seleziona gli atomi che viaggiano entro un cono stretto attorno a una direzione, che indichiamo con $\hat{y}$;
    \item il \textbf{magnete}, lungo $L$, che genera il campo disomogeneo e nel quale gli atomi vengono deflessi;
    \item una \textbf{lastra rivelatrice}, posta a distanza $D$ dall'uscita del magnete, su cui gli atomi si depositano.
\end{enumerate}

La geometria dei poli è l'elemento cruciale. Un polo termina a \emph{lama}, cioè con uno spigolo affilato parallelo alla direzione del fascio; il polo opposto presenta una \emph{scanalatura} concava affacciata allo spigolo. 
In questo modo le linee di campo si addensano in prossimità della lama e si diradano verso la scanalatura, producendo nella regione intermedia un gradiente $\partial B_z / \partial z$ molto intenso. Il fascio viene fatto passare proprio in tale regione, parallelamente allo spigolo, così da attraversare una zona in cui il gradiente è approssimativamente uniforme.

\begin{figure}[htbp]
\centering
\begin{tikzpicture}[
    line join=round, line cap=round,
    every node/.style={font=\footnotesize},
    piece/.style={line width=0.7pt},
    beam/.style={dashed, dash pattern=on 2.6pt off 2.2pt, line width=0.55pt},
    ax/.style={-{Latex[length=3.6pt]}, line width=0.45pt},
    fld/.style={-{Latex[length=3.2pt]}, line width=0.35pt},
]
 
%% ============================= (a) apparato =============================
 
% forno
\draw[piece] (0,-0.5) rectangle (1.2,0.5);
\draw[piece] (1.2,0.5) -- (1.45,0.16) (1.2,-0.5) -- (1.45,-0.16);
\node[anchor=north] at (0.6,-0.68) {forno};
 
% collimatori
\foreach \x in {2.15,2.85}{
  \draw[line width=1.1pt] (\x,0.1) -- (\x,0.46);
  \draw[line width=1.1pt] (\x,-0.1) -- (\x,-0.46);
}
\node[anchor=north] at (2.5,-0.64) {collimatori};
 
% fascio incidente
\draw[beam] (1.45,0) -- (4.1,0);
 
% atomo con momento magnetico intrinseco
\fill (3.42,0) circle (1.5pt);
\draw[-{Latex[length=3.6pt]}, line width=0.65pt] (3.42,0) -- ++(66:0.58)
      node[anchor=south, inner sep=1.5pt] {$\boldsymbol\mu$};
\draw[-{Latex[length=3.6pt]}, line width=0.65pt] (3.42,0) -- ++(0:0.55)
      node[anchor=north, inner sep=2pt] {$v$};
 
% magnete, vista laterale
\draw[piece] (4.1,0.32) rectangle (6.0,0.95);
\draw[piece] (4.1,-0.32) rectangle (6.0,-0.95);
\node at (5.05,0.635) {S};
\node at (5.05,-0.635) {N};
\node[anchor=north] at (5.05,-1.1) {magnete};
 
% lunghezza del magnete
\draw[{Latex[length=3pt]}-{Latex[length=3pt]}, line width=0.4pt]
   (4.1,1.18) -- (6.0,1.18) node[midway, anchor=south, inner sep=1.5pt] {$L$};
 
% traiettorie deflesse
\draw[line width=0.55pt] (4.1,0) .. controls (5.3,0) and (5.6,0.05) .. (6.0,0.16)
      -- (7.95,0.44);
\draw[line width=0.55pt] (4.1,0) .. controls (5.3,0) and (5.6,-0.05) .. (6.0,-0.16)
      -- (7.95,-0.44);
 
% lastra rivelatrice e tracce
\draw[piece] (7.95,-1.0) -- (7.95,1.0);
\fill (7.95,0.44) circle (1.9pt);
\fill (7.95,-0.44) circle (1.9pt);
\node[anchor=north] at (7.95,-1.15) {lastra};
 
% banda attesa classicamente
\draw[line width=0.45pt] (8.18,0.62) -- (8.32,0.62) -- (8.32,-0.62) -- (8.18,-0.62);
\node[rotate=-90, anchor=south, font=\scriptsize] at (8.38,0) {previsione classica};
 
% assi
\begin{scope}[shift={(0.45,1.25)}]
  \draw[ax] (0,0) -- (0,0.45)      node[anchor=south,      inner sep=1.5pt] {$z$};
  \draw[ax] (0,0) -- (0.45,0)      node[anchor=west,       inner sep=1.5pt] {$y$};
  \draw[ax] (0,0) -- (-0.27,-0.27) node[anchor=north east, inner sep=0.5pt] {$x$};
\end{scope}
 
\node[anchor=north west, font=\footnotesize\itshape] at (-0.15,-1.45) {(a)};
 
%% ========================== (b) sezione dei poli ==========================
\begin{scope}[shift={(10.45,0)}]
 
  % polo superiore scanalato
  \draw[piece] (-0.85,0.34) -- (-0.85,1.05) -- (0.85,1.05) -- (0.85,0.34)
     .. controls (0.55,0.34) and (0.38,0.68) .. (0,0.68)
     .. controls (-0.38,0.68) and (-0.55,0.34) .. (-0.85,0.34) -- cycle;
  \node at (0,0.87) {S};
 
  % polo inferiore a lama
  \draw[piece] (-0.85,-1.05) -- (0.85,-1.05) -- (0.85,-0.52) -- (0.26,-0.52)
     -- (0,-0.18) -- (-0.26,-0.52) -- (-0.85,-0.52) -- cycle;
  \node at (0,-0.84) {N};
 
  % linee di campo
  \draw[fld] (0,-0.14) -- (0,0.60);
  \foreach \s in {-1,1}{
    \draw[fld] (\s*0.04,-0.15) .. controls (\s*0.26,0.18) .. (\s*0.28,0.56);
    \draw[fld] (\s*0.08,-0.16) .. controls (\s*0.50,0.16) .. (\s*0.56,0.40);
  }
 
  % gradiente
  \draw[-{Latex[length=3.6pt]}, line width=0.5pt] (1.2,-0.3) -- (1.2,0.3);
  \node[anchor=west, font=\scriptsize] at (1.3,0) {$\dfrac{\partial B_z}{\partial z}$};
 
  % assi
  \draw[ax] (-1.42,-0.22) -- (-1.42,0.23) node[anchor=south, inner sep=1.5pt] {$z$};
  \draw[ax] (-1.42,-0.22) -- (-1.02,-0.22) node[anchor=north, inner sep=1.5pt] {$x$};
 
  \node[anchor=north west, font=\footnotesize\itshape] at (-1.5,-1.45) {(b)};
\end{scope}
 
\end{tikzpicture}
\caption[Apparato di Stern--Gerlach]{%
\textbf{(a)} Schema dell'apparato. 
\textbf{(b)} Sezione dei poli.}
\label{fig:stern-gerlach}
\end{figure}

Andremo a chiamare questo strumento stern-gerlach\mkeyword{stern-gerlach}.

Se le particelle fossero classiche e rotanti su loro stesse ("dotate di spin"), allora la distribuzione del loro momento angolare di spin sarebbe casuale e ogni particella sarebbe deflessa verso l'alto o il basso (verso valori positivi o negativi di z) di una differente quantità, producendo una distribuzione regolare sullo schermo di un rivelatore.

Sperimentalmente, però, \emph{le particelle che passano attraverso il campo magnetico sono deflesse o in alto o in basso di una quantità precisa}.\\
Inoltre, se ruotiamo di un certo angolo lo stern-gerlach otteniamo lo stesso risultato solo che i punti di delessione adesso sono paralleli all'asse ruotato. 


\begin{figure}[htbp]
\centering
\begin{tikzpicture}[
    line join=round, line cap=round,
    every node/.style={font=\footnotesize},
    piece/.style={line width=0.7pt},
    beam/.style={dashed, dash pattern=on 2.6pt off 2.2pt, line width=0.5pt},
    traj/.style={line width=0.4pt},
    curve/.style={line width=0.8pt},
    ax/.style={-{Latex[length=3.6pt]}, line width=0.45pt},
]
 
%% ######################## (a) previsione classica ########################
\begin{scope}[shift={(0,3.9)}]
 
  % fascio incidente
  \draw[beam] (-1.5,0) -- (0,0);
 
  % magnete
  \draw[piece] (0,0.34) rectangle (1.9,0.95);
  \draw[piece] (0,-0.34) rectangle (1.9,-0.95);
  \node at (0.95,0.645) {S};
  \node at (0.95,-0.645) {N};
 
  % ventaglio di traiettorie
  \foreach \z in {-1.25,-0.94,-0.63,-0.31,0,0.31,0.63,0.94,1.25}{
    \draw[traj] (0,0) .. controls (1.15,0) and (1.5,{0.07*\z})
                .. (1.9,{0.176*\z}) -- (4.45,\z);
  }
 
  % lastra
  \draw[piece] (4.45,-1.75) -- (4.45,1.75);
 
  % distribuzione continua
  \draw[curve] plot[domain=-1.72:1.72, samples=140, smooth]
      ({4.68+1.25*exp(-(\x*\x)/0.5408)}, {\x});
 
  % etichette
  \node[anchor=west, align=left] at (6.0,0.95)
     {distribuzione\\ continua};
  \draw[line width=0.4pt] (5.95,0.95) -- (4.95,0.78);
 
  \node[anchor=east, align=right, font=\footnotesize] at (-1.65,0)
     {\textbf{(a)} previsione\\ classica};
 
  % assi
  \draw[ax] (-0.95,1.25) -- (-0.95,1.72) node[anchor=south, inner sep=1.5pt] {$z$};
  \draw[ax] (-0.95,1.25) -- (-0.48,1.25) node[anchor=west,  inner sep=1.5pt] {$y$};
\end{scope}
 
%% ######################## (b) risultato osservato ########################
\begin{scope}[shift={(0,0)}]
 
  % fascio incidente
  \draw[beam] (-1.5,0) -- (0,0);
 
  % magnete
  \draw[piece] (0,0.34) rectangle (1.9,0.95);
  \draw[piece] (0,-0.34) rectangle (1.9,-0.95);
  \node at (0.95,0.645) {S};
  \node at (0.95,-0.645) {N};
 
  % due sole traiettorie
  \draw[line width=0.55pt] (0,0) .. controls (1.15,0) and (1.5,0.05)
       .. (1.9,0.127) -- (4.45,0.72);
  \draw[line width=0.55pt] (0,0) .. controls (1.15,0) and (1.5,-0.05)
       .. (1.9,-0.127) -- (4.45,-0.72);
 
  % lastra
  \draw[piece] (4.45,-1.75) -- (4.45,1.75);
 
  % due picchi stretti
  \draw[curve] plot[domain=-1.72:1.72, samples=320, smooth]
      ({4.68+1.25*exp(-((\x-0.72)*(\x-0.72))/0.0242)
            +1.25*exp(-((\x+0.72)*(\x+0.72))/0.0242)}, {\x});
 
  % separazione
  \draw[{Latex[length=3pt]}-{Latex[length=3pt]}, line width=0.4pt]
     (4.2,-0.72) -- (4.2,0.72) node[midway, fill=white, inner sep=1pt]
     {$2\Delta z$};
 
  % etichette dei picchi
  \node[anchor=west] at (6.0,0.78)  {$S_z = +\hbar/2$};
  \node[anchor=west] at (6.0,-0.78) {$S_z = -\hbar/2$};
 
  \node[anchor=east, align=right, font=\footnotesize] at (-1.65,0)
     {\textbf{(b)} risultato\\ osservato};
\end{scope}
 
\end{tikzpicture}
\caption[Stern--Gerlach: previsione classica e risultato]{%
\textbf{(a)} Previsione classica.
\textbf{(b)} Risultato osservato.}
\label{fig:sg-distribuzioni}
\end{figure}

\newpage

Continuiamo gli esperimenti allineando più stern gerlach. 
 
Per alleggerire gli schemi conviene adottare una notazione compatta: indichiamo ciascun apparato con una scatola contrassegnata dall'asse lungo cui è orientato il gradiente, $\hat{z}$, $\hat{x}$ e così via. 
Il fascio entra da sinistra ed esce dai due canali di deflessione, che disegniamo come due rami. Un tratto spesso su un ramo indica che quel fascio viene \emph{bloccato}, cioè intercettato da uno schermo assorbente e rimosso dall'esperimento. Le percentuali riportate sui rami sono le frazioni di atomi raccolte in uscita.
 
 
\subsection{Creazione di una sorgente pura}
 
Il fascio che esce dal forno è casuale: non siamo in grado di prevedere se un singolo atomo verrà deflesso verso l'alto o verso il basso, e le due uscite raccolgono ciascuna il $50\%$ degli atomi. 
Se riallineiamo i due fasci uscenti e li rimandiamo in un secondo apparato, qualunque sia la sua orientazione ritroviamo di nuovo una ripartizione $50\%$--$50\%$: il fascio ricostruito si comporta esattamente come quello di partenza. 
Chiamiamo sorgente mista\mkeyword{sorgente mista} una sorgente di questo tipo.
 
Blocchiamo invece una delle due uscite, per esempio quella inferiore, e teniamo il solo fascio deflesso verso l'alto. Se mandiamo questo fascio in un secondo apparato orientato anch'esso lungo $\hat{z}$, osserviamo che
%
\begin{itemize}
    \item il $100\%$ degli atomi esce dal canale superiore;
    \item il canale inferiore non raccoglie alcun atomo.
\end{itemize}
%
Il risultato è riproducibile: per quanti apparati orientati lungo $\hat{z}$ mettiamo in fila, l'esito resta sempre lo stesso e nessun atomo viene perso.
Chiamiamo il dispositivo così costruito sorgente pura\mkeyword{sorgente pura} lungo $\hat{z}\!\uparrow$, e lo indichiamo con il simbolo compatto di Fig.~\ref{fig:sorgente-pura}. Analogamente si costruisce una sorgente pura lungo qualunque altra direzione, bloccando il canale opportuno di un apparato orientato lungo quella direzione.
 
\begin{figure}[htbp]
\centering
\begin{tikzpicture}[sgstyle]
 
% --- (a) sorgente non polarizzata
\begin{scope}[shift={(0,1.9)}]
  \forno{0,-0.31}
  \draw[ray] (0.95,0) -- (1.9,0);
  \node[sg] at (2.31,0) {$\hat z$};
  \draw[ray] (2.72,0) -- (3.9,0.5);
  \draw[ray] (2.72,0) -- (3.9,-0.5);
  \node[lbl, anchor=west] at (3.95,0.5)  {$50\%$};
  \node[lbl, anchor=west] at (3.95,-0.5) {$50\%$};
  \node[anchor=east, font=\footnotesize] at (-0.25,0) {(a)};
\end{scope}
 
% --- (b) uscita inferiore bloccata
\begin{scope}[shift={(0,-0.6)}]
  \forno{0,-0.31}
  \draw[ray] (0.95,0) -- (1.9,0);
  \node[sg] at (2.31,0) {$\hat z$};
  \draw[ray] (2.72,0) -- (4.3,0.5);
  \draw[ray] (2.72,0) -- (3.75,-0.42);
  \draw[stop] (3.82,-0.62) -- (3.82,-0.26);
  \node[lbl, anchor=west] at (4.35,0.5) {$100\%$};
  \node[anchor=east, font=\footnotesize] at (-0.25,0) {(b)};
 
  \node[font=\footnotesize] at (5.65,0) {$\equiv$};
  \sorgentepura{6.1,-0.34}{$\hat z\!\uparrow$}
  \draw[ray] (7.08,0) -- (7.8,0);
\end{scope}
 
\end{tikzpicture}
\caption[Costruzione di una sorgente pura]{%
\textbf{(a)} Il fascio uscente dal forno si ripartisce equamente fra i due
canali. \textbf{(b)} Bloccando il canale inferiore si ottiene un fascio che,
rianalizzato lungo $\hat{z}$, esce interamente dal canale superiore: è la
sorgente pura lungo $\hat{z}\!\uparrow$, per la quale adottiamo il simbolo
compatto a destra.}
\label{fig:sorgente-pura}
\end{figure}
 
 
\subsection{La sequenza $z \to x \to z$}
 
Prendiamo una sorgente pura lungo $\hat{z}\!\uparrow$ e analizziamola con apparati orientati lungo assi diversi.
 
\begin{enumerate}
    \item Con un apparato lungo $\hat{z}$ ritroviamo quanto appena visto: il $100\%$ degli atomi esce dal canale superiore.
 
    \item Con un apparato lungo $\hat{x}$ il fascio si ripartisce equamente, nella proporzione $50\%$--$50\%$: la selezione operata lungo $\hat{z}$ non privilegia in alcun modo uno dei due canali lungo $\hat{x}$.
 
    \item Blocchiamo ora il canale $\hat{x}\!\downarrow$ e mandiamo il solo fascio $\hat{x}\!\uparrow$ in un terzo apparato, orientato di nuovo lungo $\hat{z}$. Otteniamo ancora una ripartizione $50\%$--$50\%$, cioè il $25\%$ e il $25\%$ degli atomi entrati nel secondo apparato.
\end{enumerate}
 
Il terzo risultato va confrontato con il primo. Se al posto dell'apparato lungo $\hat{x}$ ne inseriamo uno lungo $\hat{z}$, il fascio prosegue indisturbato e l'apparato finale raccoglie il $100\%$ degli atomi nel canale superiore; inserendo invece l'apparato lungo $\hat{x}$ e bloccandone un canale, l'apparato finale raccoglie metà degli atomi nel canale superiore e metà in quello inferiore.
 
\begin{figure}[htbp]
\centering
\begin{tikzpicture}[sgstyle]
 
% --- (a)
\begin{scope}[shift={(0,2.5)}]
  \sorgentepura{0,-0.34}{$\hat z\!\uparrow$}
  \draw[ray] (0.98,0) -- (1.75,0);
  \node[sg] at (2.16,0) {$\hat z$};
  \draw[ray] (2.57,0) -- (3.6,0.45);
  \draw[ray] (2.57,0) -- (3.6,-0.45);
  \node[lbl, anchor=west] at (3.65,0.45)  {$100\%$};
  \node[lbl, anchor=west] at (3.65,-0.45) {$0\%$};
  \node[anchor=east, font=\footnotesize] at (-0.25,0) {(a)};
\end{scope}
 
% --- (b)
\begin{scope}[shift={(0,0.9)}]
  \sorgentepura{0,-0.34}{$\hat z\!\uparrow$}
  \draw[ray] (0.98,0) -- (1.75,0);
  \node[sg] at (2.16,0) {$\hat x$};
  \draw[ray] (2.57,0) -- (3.6,0.45);
  \draw[ray] (2.57,0) -- (3.6,-0.45);
  \node[lbl, anchor=west] at (3.65,0.45)  {$50\%$};
  \node[lbl, anchor=west] at (3.65,-0.45) {$50\%$};
  \node[anchor=east, font=\footnotesize] at (-0.25,0) {(b)};
\end{scope}
 
% --- (c)
\begin{scope}[shift={(0,-1.1)}]
  \sorgentepura{0,-0.34}{$\hat z\!\uparrow$}
  \draw[ray] (0.98,0) -- (1.75,0);
  \node[sg] at (2.16,0) {$\hat x$};
  \draw[ray] (2.57,0) -- (3.55,0.5);
  \draw[ray] (2.57,0) -- (3.35,-0.4);
  \draw[stop] (3.42,-0.58) -- (3.42,-0.24);
  \node[lbl, anchor=south] at (3.45,0.52) {$50\%$};
  \draw[ray] (3.55,0.5) -- (4.25,0.5);
  \node[sg] at (4.66,0.5) {$\hat z$};
  \draw[ray] (5.07,0.5) -- (6.0,0.92);
  \draw[ray] (5.07,0.5) -- (6.0,0.08);
  \node[lbl, anchor=west] at (6.05,0.92) {$25\%$};
  \node[lbl, anchor=west] at (6.05,0.08) {$25\%$};
  \node[anchor=east, font=\footnotesize] at (-0.25,0) {(c)};
\end{scope}
 
\end{tikzpicture}
\caption[Sequenza $z \to x \to z$]{%
Analisi di una sorgente pura lungo $\hat{z}\!\uparrow$.
\textbf{(a)} Un apparato lungo $\hat{z}$ lascia passare tutti gli atomi dal
canale superiore. \textbf{(b)} Un apparato lungo $\hat{x}$ ripartisce il fascio
equamente. \textbf{(c)} Selezionando $\hat{x}\!\uparrow$ e rianalizzando lungo
$\hat{z}$, le percentuali sono riferite al fascio entrante nel primo apparato.}
\label{fig:sequenza-zxz}
\end{figure}
 
 
\subsection{Ricombinazione dei fasci}
 
Negli esperimenti precedenti uno dei due canali veniva sempre bloccato.
Costruiamo ora un dispositivo che li conservi entrambi: un apparato orientato lungo $\hat{x}$ separa il fascio nei due canali, che vengono fatti proseguire per un tratto senza essere né intercettati né rivelati, e un secondo apparato identico al primo ma di gradiente invertito li riporta a sovrapporsi in un unico fascio. 
Il fascio ricombinato viene infine analizzato con un apparato lungo $\hat{z}$.
 
Il risultato è che il $100\%$ degli atomi esce dal canale superiore e nessuno da quello inferiore, \emph{esattamente come se i due apparati lungo $\hat{x}$ non ci fossero}. 
In particolare non si osserva la ripartizione $50\%$--$50\%$ che si otteneva bloccando uno dei due canali.
 
\subsection{Ricombinazione dei fasci con osservatore}
 
Ripetiamo l'esperimento precedente aggiungendo, lungo uno dei due cammini separati, un dispositivo in grado di rivelare il passaggio dell'atomo e di registrare quale dei due cammini esso abbia percorso; chiamiamolo osservatore\mkeyword{osservatore}. Nessuno dei due fasci viene bloccato, e l'intensità complessiva in uscita non cambia.
 
L'apparato finale lungo $\hat{z}$ raccoglie ora $50\%$ degli atomi nel canale superiore e $50\%$ in quello inferiore. Vale anche l'implicazione inversa: se osserviamo una ripartizione $50\%$--$50\%$ possiamo concludere che lungo i cammini era presente un osservatore.
 
Due osservazioni completano il quadro.
 
\begin{itemize}
    \item Ciò che conta è che il dispositivo \emph{lasci traccia} del cammino percorso. Un apparato che interagisce con gli atomi senza registrare alcuna informazione su quale cammino sia stato seguito non altera il risultato, e si continua a raccogliere il $100\%$ degli atomi nel canale superiore.
 
    \item Non è la misura in sé a modificare l'esito, ma il fatto che essa distingua i due cammini. Se collochiamo lo stesso rivelatore di posizione \emph{dopo} la ricombinazione, cioè fra il secondo apparato lungo $\hat{x}$ e quello lungo $\hat{z}$, dove esiste ormai un solo fascio, l'esito torna a essere $100\%$ nel canale superiore.
\end{itemize}
 
\begin{figure}[htbp]
\centering
\begin{tikzpicture}[sgstyle]
 
% --- (a) senza osservatore
\begin{scope}[shift={(0,2.1)}]
  \sorgentepura{0,-0.34}{$\hat z\!\uparrow$}
  \draw[ray] (0.98,0) -- (1.6,0);
  \node[sg] at (2.01,0) {$\hat x$};
  \draw[ray] (2.42,0) .. controls (2.9,0) and (3.0,0.62) .. (3.5,0.62)
             -- (4.5,0.62) .. controls (5.0,0.62) and (5.1,0) .. (5.58,0);
  \draw[ray] (2.42,0) .. controls (2.9,0) and (3.0,-0.62) .. (3.5,-0.62)
             -- (4.5,-0.62) .. controls (5.0,-0.62) and (5.1,0) .. (5.58,0);
  \node[sg] at (5.99,0) {$\hat x$};
  \draw[ray] (6.4,0) -- (6.95,0);
  \node[sg] at (7.36,0) {$\hat z$};
  \draw[ray] (7.77,0) -- (8.65,0.42);
  \draw[ray] (7.77,0) -- (8.65,-0.42);
  \node[lbl, anchor=west] at (8.7,0.42)  {$100\%$};
  \node[lbl, anchor=west] at (8.7,-0.42) {$0\%$};
  \node[lbl, anchor=north] at (2.01,-0.45) {separa};
  \node[lbl, anchor=north] at (5.99,-0.45) {ricombina};
  \node[anchor=east, font=\footnotesize] at (-0.25,0) {(a)};
\end{scope}
 
% --- (b) con osservatore
\begin{scope}[shift={(0,-1.4)}]
  \sorgentepura{0,-0.34}{$\hat z\!\uparrow$}
  \draw[ray] (0.98,0) -- (1.6,0);
  \node[sg] at (2.01,0) {$\hat x$};
  \draw[ray] (2.42,0) .. controls (2.9,0) and (3.0,0.62) .. (3.5,0.62)
             -- (4.5,0.62) .. controls (5.0,0.62) and (5.1,0) .. (5.58,0);
  \draw[ray] (2.42,0) .. controls (2.9,0) and (3.0,-0.62) .. (3.5,-0.62)
             -- (4.5,-0.62) .. controls (5.0,-0.62) and (5.1,0) .. (5.58,0);
  \draw[line width=0.7pt] (3.78,0.38) rectangle (4.22,0.86);
  \node[lbl, anchor=south] at (4.0,0.95) {osservatore};
  \node[sg] at (5.99,0) {$\hat x$};
  \draw[ray] (6.4,0) -- (6.95,0);
  \node[sg] at (7.36,0) {$\hat z$};
  \draw[ray] (7.77,0) -- (8.65,0.42);
  \draw[ray] (7.77,0) -- (8.65,-0.42);
  \node[lbl, anchor=west] at (8.7,0.42)  {$50\%$};
  \node[lbl, anchor=west] at (8.7,-0.42) {$50\%$};
  \node[anchor=east, font=\footnotesize] at (-0.25,0) {(b)};
\end{scope}
 
\end{tikzpicture}
\caption[Ricombinazione dei fasci]{%
\textbf{(a)} I due canali non vengono bloccati ma ricombinati da un secondo
apparato lungo $\hat{x}$: l'analisi finale lungo $\hat{z}$ restituisce il
$100\%$ degli atomi nel canale superiore, come se i due apparati non ci fossero.
\textbf{(b)} Inserendo lungo uno dei cammini un dispositivo che registri quale
di essi sia stato percorso, l'analisi finale dà $50\%$--$50\%$.}
\label{fig:ricombinazione}
\end{figure}

\section{Interpretazione dei risultati}

Riportiamo in forma compatta ciò che abbiamo analizzato nei diversi esperimenti. Ogni stern-gerlach è orientato in un direzione $\hat{n}$ a scelta e chiamiamo le due uscite $+$ e $-$. 

\begin{enumerate}
    \item Qualunque sia $\hat{n}$, le uscite sono sempre due e la separazione ha sempre la stessa ampiezza. 
    \item Prendendo l'uscita $+$ e rimandandola in un altro filtro $\hat{n}$ otteniamo un altro fascio $+$. 
    \item Con il fascio appena uscito dal forno, qualunque $\hat n$ dà $50/50$. Non esiste alcuna direzione privilegiata nello spazio.
    \item Preparando un filtro $\hat{n}$, prendendo l'uscita $+$ e mandandola in un filtro $\hat{m}$ con angolo $\theta$ fra i due assi si misura 
    \begin{equation}
      \frac{I_+}{I_{\text{tot}}} = \cos^2 \frac{\theta}{2}
    \end{equation}
    In particolare per $\theta = 90^{\circ}$ si ottiene che il rapporto è $1/2$.
    \item I filtri non commutano.
    \item Lo sdoppiamento è reversibile se lungo l'apparato di ricollimazione non si osserva nulla. 
    \item L’informazione sul cammino distrugge tutto. Se si considera lo stesso dispositivo ma si mette un rivelatore che registri quale cammino è stato percorso, senza bloccare nulla, il filtro $z$ finale dà $50/50$, a intensità piena.
\end{enumerate}

\subsection{Conclusioni} 

Dai dati osservati allora si possono fare le seguenti conclusioni. 

\begin{enumerate}[label =\Roman*]
  \item \emph{Esiste una proprietà a due valori e vive sulle direzioni}. Da 1) e 3) si conclude che non c’è una “quantità orientata” con un suo asse, perché ogni direzione che si sceglie si comporta allo stesso modo.
  \item \emph{Domande lungo direzioni diverse non sono simultaneamente ponibili}. Dalla 5) si sa che se il sistema portasse con sé una risposta per ogni direzione, fissata e indisturbata, il terzo filtro dovrebbe ridare $100\%$.
  \item \emph{Uno stato preparato è individuato da una direzione}. 
  \item \emph{Due stati sono perfettamente distinguibili quando le direzioni sono opposte}. Da 4) con $\theta=180^{\circ}$: probabilità zero.
  \item \emph{Lo sdoppiamento non è una suddivisione della popolazione}. Da 6) si ha che se dopo il magnete $z$ il fascio fosse “metà atomi del tipo $+z$, metà del tipo $-z$”, allora riunirli rimetterebbe insieme due popolazioni, e l’informazione su $x$ — che quelle due popolazioni, prese separatamente, non hanno — resterebbe persa. Invece torna intatta.
  \item \emph{La distruzione non è dovuta alla separazione ma all'informazione}. Nel punto 6) e 7) abbiamo lo stesso dispositivo, stessa intensità, stessa meccanica: cambia solo che esiste da qualche parte un registro di quale cammino. Questo basta a distruggere l'informazione. 
\end{enumerate}

\subsection{Inferenza della necessità degli assiomi}

A partire dalle concluisioni ottenute è chiaro che la meccanica classica non è sufficiente a spiegare questi fenomeni, è allora necessario svillupare un nuovo strumento, la Meccanica Quantistica.

Per andare a porre le basi di questo nuovo strumento sfruttiamo le conclusioni dell'esperimento. 

\subsection*{Struttura lineare}

Si deve rappresentare uno stato con un oggetto che si spezza in due parti e si ricompone esattamente. Non “si divide in due sottopopolazioni”, ma proprio si decompone e si risomma.

L’unica struttura matematica in cui decomporre e ricomporre è esatto e senza perdita è uno spazio vettoriale: lo stato è una somma di componenti lungo i due cammini, e la ricombinazione è letteralmente l’addizione. 

\subsection*{Numeri complessi}

Siano $(c_+,c_-)$ i due numeri associati alle due uscite del filtro. Resta da stabilire di che tipo di numeri si tratti, e a rispondere basta un conteggio di parametri.

Due vincoli sono già fissati: la normalizzazione, perché nessun atomo viene perso, e l'irrilevanza di un fattore moltiplicativo comune a entrambe le ampiezze, che moltiplicherebbe tutte le intensità per lo stesso numero senza alcun effetto osservabile.

Se le ampiezze fossero reali avremmo due numeri meno la normalizzazione, cioè un solo parametro libero: una circonferenza. Il fattore globale, in questo caso un semplice segno, non ne toglie alcuno, essendo un'identificazione discreta e non una famiglia continua. Ma dalla conclusione III gli stati preparabili sono in corrispondenza con le direzioni dello spazio, cioè con una \emph{sfera}, che di parametri ne richiede due. Con ampiezze reali non riusciremmo a distinguere uno stato preparato lungo $\hat{y}$ da uno preparato lungo $\hat{x}$, mentre sperimentalmente li distinguiamo senza difficoltà.

Se le ampiezze sono complesse abbiamo invece quattro numeri reali meno la normalizzazione, e in più il fattore globale è ora una fase $e^{i\alpha}$, cioè una famiglia continua che toglie un ulteriore parametro: ne restano esattamente due.

Lo spazio degli stati di un singolo filtro è dunque $\mathbb{C}^2$.

\subsection*{Prodotto scalare}

Serve una regola che da $(c_+, c_-)$ si deducano le probabilità. \\ 
Inoltre la somma deve fare 1 e restare 1 sotto le operazioni reversibili di 6), perché non si perdono atomi.

La forma quadratica definita positiva che soddisfa entrambi è $|c_\pm|^2$, da cui $|c_+|^2+|c_-|^2=1$. Questo introduce il prodotto scalare hermitiano, e con esso la norma.

\subsection*{Operatori Hermitiani}

Che cosa è, sperimentalmente, “il filtro $\hat{n}$”:
\begin{itemize}
  \item due esiti, etichettati da due numeri reali (ciò che leggi: la deflessione, e quindi $\pm\hbar/2$);
  \item due stati che danno quei due esiti con certezza;
  \item quei due stati sono perfettamente distinguibili, cioè uno dà probabilità zero per l’esito dell’altro, e quindi hanno prodotto scalare nullo: sono ortogonali;
  \item ogni altro stato si decompone su quei due, quindi formano una base.
\end{itemize}

Quindi un filtro è esattamente: una base ortonormale di $\mathbb C^2$, più due numeri reali. L’oggetto matematico che contiene precisamente questa informazione, e nient’altro, è

$$\hat A = \sum_i \lambda_i\,|e_i\rangle\langle e_i|,\qquad \lambda_i\in\mathbb R,\quad \{|e_i\rangle\}\ \text{ortonormale}$$

cioè, per il teorema spettrale, un operatore hermitiano. 

Un altro fattore che suggerisce la natura da operatori dei filtri è la non commutatività.%\footnote{Fino ad ora abbiamo spiegato la necessità degli elementi nell'assioma 1 e assioma 2, tuttavia la spiegazione per l'assioma 3 è molto meno banale e per cui si consiglia di guardare \emph{Modern Quantum Mechanic} di J.J. Sakurai.}

\subsection*{Valori medi}

Gli assiomi 1 e 2 fissano lo spazio degli stati e i valori che una misura può
restituire, ma non collegano ancora lo stato alle frequenze osservate. La regola
mancante l'abbiamo però già ricavata: detta $\{\ket{e_n}\}$ la base associata al
filtro, le ampiezze sono $c_n = \braket{e_n|\psi}$ e la sezione precedente ha
mostrato che le probabilità sono $P_n = |\braket{e_n|\psi}|^2$. Il valor medio è
allora la definizione ordinaria di valore atteso, e usando la rappresentazione
spettrale dell'assioma 2
%
\begin{equation}
    \sum_n \lambda_n P_n
    = \sum_n \lambda_n \braket{\psi|e_n}\braket{e_n|\psi}
    = \Braket{\psi|\hat{O}|\psi}
\end{equation}
%
Conviene enunciare l'assioma sui valori medi perché questo contiene già le
probabilità: applicandolo al proiettore $\hat{P}_n = \ket{e_n}\bra{e_n}$ si
ottiene $\langle P_n \rangle = P_n$.\footnote{A differenza dei primi due, questo
assioma non è imposto dagli esperimenti, che forniscono le sole frequenze. Il
teorema di Gleason (1957) mostra però che in dimensione maggiore o uguale a $3$
esso non è indipendente dagli altri due.}